
\begin{frame}{Economist 2016 Election Model \citep{economist:2020:election}}

\begin{minipage}[t]{0.35\textwidth}
    \ElectionData{}
\end{minipage}
\begin{minipage}{0.04\textwidth}
\hfill
\end{minipage}
\begin{minipage}[t]{0.59\textwidth}

A time series model to predict the 2016 US presidential election
outcome from polling data.

%\pause
\spskip
\begin{itemize}
\item $\xvec = \x_1, \ldots, \x_N =$ polling data.
\item $\theta = $ parameters for everything we don't know:
%
\begin{itemize}
\item Daily information shocks (time series residuals)
\item Ideosyncracies of particular polling agencies
\item Biases in polling methods
\end{itemize}
%
\item Model polling randomness with $\p(\xvec | \theta)$
\item $f(\theta) = $ Democratic \% of vote on election day
\end{itemize}
%
\end{minipage}

\pause

\question{Many values of the unknowns $\theta$ are consistent with the data.\\
How can we account for this uncertainty when making election predictions?}
\pause
\textbf{(Bayesian) Answer: }
    Average over values that are plausibly consistent with the data.

\pause
\textbf{Formally:} Define a prior $\p(\theta)$, 
compute the posterior $\p(\theta \vert \xvec)$, and estimate
%
\begin{itemize}
\item[] $f(\theta) \approx \expect{p(\theta \vert \xvec)}{f(\theta)}$, with 
uncertainty proportional to $\sqrt{\var{p(\theta \vert \xvec)}{f(\theta)}}$.
\end{itemize}
\pause
% \spskip
\textbf{Practically:} We compute Markov chain Monte Carlo (MCMC) draws from
$p(\theta \vert \xvec)$.


\end{frame}



%%%%%%%%%%%%%%%%%%%%%%%%%%%%%%%%%
%%%%%%%%%%%%%%%%%%%%%%%%%%%%%%%%%
%%%%%%%%%%%%%%%%%%%%%%%%%%%%%%%%%



\begin{frame}{Economist 2016 Election Model \citep{economist:2020:election}}

\begin{minipage}[t]{0.35\textwidth}
    \ElectionData{}
\end{minipage}
\begin{minipage}{0.04\textwidth}
\hfill
\end{minipage}
\begin{minipage}[t]{0.59\textwidth}

A time series model to predict the 2016 US presidential election
outcome from polling data.

%\pause
\spskip
\begin{itemize}
\item $\xvec = \x_1, \ldots, \x_N =$ Polling data ($N = 361$).
\item $\theta = $ Lots of random effects (day, pollster, etc.)
\item $f(\theta) = $ Democratic \% of vote on election day
\end{itemize}

%\pause
\spskip
We want to know $\expect{p(\theta \vert \xvec)}{f(\theta)}$.

%\pause
\spskip
Typically, we compute Markov chain Monte Carlo (MCMC) draws from the
posterior $p(\theta \vert \xvec)$.


\end{minipage}


\pause
\question{
    The people who responded to the polls were randomly selected.  \\
    If we had selected a different random sample, how much would our estimate have changed?\\
    How can we estimate $\var{\xfdist}{\expect{p(\theta \vert \xvec)}{f(\theta)}}$?
    }

\pause
%
\begin{itemize}
\item Except in special cases, we expect 
    $\var{p(\theta \vert \xvec)}{f(\theta)} \ne \var{\xfdist}{\expect{p(\theta \vert \xvec)}{f(\theta)}}$
% \item The model is (surely) misspecified
% \item The posterior expectation marginalizes over many nuisance parameters
\item We are interested re--sampling for \textit{this} election, not a hypothetical future election
\end{itemize}
%

% }

\end{frame}

    
%%%%%%%%%%%%%%%%%%%%%%%%%%%%%%%%%
%%%%%%%%%%%%%%%%%%%%%%%%%%%%%%%%%
%%%%%%%%%%%%%%%%%%%%%%%%%%%%%%%%%




\begin{frame}{Results}

\question{
    \textbf{Idea: } Re-fit with bootstrap samples of data \citep{huggins:2023:bayesbag}\\
    \textbf{Problem:} Each MCMC run takes about $10$ hours (Stan, six cores).\\
    \textbf{Proposal:} Use full--data posterior draws
    to form a linear approximation to {\em data reweightings}.
}

\onslide<2->{
\begin{minipage}{0.43\textwidth}
    \ElectionData{}
\end{minipage}
\begin{minipage}{0.43\textwidth}
\begin{tikzpicture}
    \node[anchor=south west,inner sep=0] (image) at (0,0) {
        \ElectionResultsGlobal{}
    };
    \onslide<4->{
    \begin{scope}[x={(image.south east)},y={(image.north west)}]
        \draw[gray, thick, <-] (0.55,0.1) -- ++(0.15,0.1)
                node[right,black,fill=white]
                {\small Both are approximations};
        \draw[gray, thick, <-] (0.55,0.3) -- ++(0.15,-0.1);
        \node[right] at (0.7, 0.15) { \small to the truth!};
    \end{scope}
    }
\end{tikzpicture}
\end{minipage}
}
%
\onslide<3->{
\vfill \begin{center}
%
\begin{center}
\begin{tabular}{ r l }
    {  Compute time for 100 bootstraps:} &
    {  51 days} \spskip\\
    {  Compute time for the linear approximation:} &
    {  Seconds}\\
    {  (But note the approximation has some error)} &
    \\
\end{tabular}
\end{center}

\end{center}
}
%
\end{frame}


